% =============================================================================
%  CHAPTER 8 -- DISCUSSION                                            [v4 OWNS]
% -----------------------------------------------------------------------------
%  v4.1 (2026-09-24): the Discussion follows the Conclusion and holds three sections
%  (author): Limitations (the v4.0 text, one scope sentence added), Towards Deployment
%  (from the v4.0 Conclusion, plus the selection-rule paragraph the author's notes asked
%  for) and Future Work (from the v4.0 Conclusion; the simulator is Isaac Sim, author).
%  The v4.0 sections Interpretation, Negative and Inconclusive Results and Threats to
%  Validity are dropped as duplicates of Chapter 6 (author) and archived in
%  withheld/20260924_v4.1_archive/07_discussion_v4.0.tex. Labels sec:disc:interpretation,
%  sec:disc:negative and sec:disc:threats are gone: nothing in Chapters 1-6 references
%  them (checked at v4.1). ch:discussion and sec:disc:limitations are kept;
%  sec:conc:practice -> sec:disc:practice, sec:conc:future -> sec:disc:future.
%  Every number restates one of Chapter 6.
%  v4.2 (2026-09-25): the ChatGPT audit of v4.1a applied (AUDIT_v4.1a_2026-09-25.md, section 15,
%  items D1-D11): opening; altitude, budget, lead/cross-track, controller-cost and candidate
%  paragraphs scoped to their populations; future-work items reworded as tests; plus the author's
%  planning-horizon item (v4.2).
%  Individual changelog: changelogs/v4.2_20260925_audit_v4.1a_applied.md
% =============================================================================

\chapter{Discussion}\label{ch:discussion}

This chapter states what the framework does not do, what a deployment has to add to it, and what has
not been tested.

\section{Limitations}\label{sec:disc:limitations}

\paragraph{The plan is a sequence of positions.} The generative model plans position increments and
never an actuator command; everything below the increment is the plant's controller
(\autoref{sec:method:setpoint}). The projection's dynamics model is a first-order integrator over the
commanded quantity, so feasibility is a statement about that model and its mismatch bound, which the
tightening absorbs, and not about the plant. On the manipulator the guarantee is conditional: the tightening
absorbs the mismatch only while the one-step change of the tracking error stays within the configured margin,
which was not calibrated here (\autoref{sec:setup:protocol:eval}); on the quadrotor whether a plan can be flown is the
controller's, which is what UAV-s-curve shows.

\paragraph{A constraint is what is declared.} The projector holds a plan inside the declared geometry,
halfspaces and keep-out regions, and protects nothing else. On UAV-pillars five of the six pillars are
in no constraint set and a third of the flights end on them. Whether a scene is safe follows from its
constraint set and not from the projection; the remedy is a constraint per obstacle, or a margin beyond
the rotor reach, and it is a decision of scene design. The span over which a wall or a plane acts is
checked at the commanded position, half a metre ahead of the vehicle on the corridor, which leaves the
hand-over residue at the end of every plane.

\paragraph{Endpoint projection is defined on a velocity field.} It adds its control to the field of a
flow-based sampler and has no counterpart for the diffusion model, whose reverse chain has none
(\autoref{sec:method:hardflow}). The comparison of the two projection methods is therefore made within
the flow-based family, and the diffusion baseline is read under per-step projection alone.

\paragraph{Plan quality is read through violating steps and the drawn plans.} The plans are scored by
success, constraint satisfaction, control steps and time, and inspected as drawn (\autoref{fig:raw-plans},
\autoref{fig:uav-scurve-plans}). No curvature, jerk or path-length measure is taken over the stored
plans, so the ribbon of the average-velocity models against the scribble of the one-step baseline, and
the rough low-budget plans that cross the s-curve most often, are read from the drawings and the counts
beside them.

\paragraph{Scope of the transfer.} The vision-conditioned task adds camera observations and contact to
the manipulator in one step; the state-observed form of the alignment task was not run, so what
D3IL-aligning shows is that the result survives a harder, differently observed task, and not the step
from state to vision alone. The quadrotor is observed through its state, at one altitude on UAV-pillars,
where the plan carries no altitude; the demonstrations of the other two scenes lie between $0.90$ and
$1.30$\,m, and the projected flights over the hump climb above that range.
D3IL-aligning, UAV-corridor and UAV-s-curve are evaluated on one training seed, D3IL-avoiding and
UAV-pillars on five (\autoref{tab:protocol}); the ten alignment contexts are drawn from the training set, and no
held-out context was evaluated. No loop of this thesis was run in real time, and no environment is a physical system.
\dataref{\autoref{sec:method:manipulator}; \autoref{sec:method:env:uav}; \autoref{sec:res:uav:pillars};
\autoref{sec:res:uav:projected}; \autoref{sec:method:hardflow}; \autoref{tab:uav-demos};
\autoref{sec:method:setpoint}.}

\section{Towards Deployment}\label{sec:disc:practice}

Three properties of the framework carry into practice as it stands. The constraints enter only at
sampling time, through the projection; nothing about them reaches training, so a new restriction is a
change of the constraint set and not a training run, and the feasible set can be replaced between two
replans, as the switched walls of the quadrotor scenes are (\autoref{sec:method:env:uav}). The step
budget is a deployment choice: a flow-based checkpoint is read at the budget selected for its task,
one sampling step on D3IL-avoiding and twenty on D3IL-aligning, where the diffusion model of
\ac{DPCC}, with the sampler inherited here, needs one checkpoint per budget.
And the plan is a sequence of positions that a tracking controller follows, which is what let one
planner fly the quadrotor unchanged.

Two quantities bound that use. A plan is held inside the declared geometry only, so the constraint set
has to name every obstacle that matters, with a margin that covers the vehicle's extent, and the
channel it leaves has to be flyable by the controller: the commanded position runs ahead of the
vehicle, by $0.49$ to $0.61$\,m at the first violating step of the projected flow-based corridor
flights and by about $0.3$\,m on the s-curve, while the unprojected flow-based flights of the s-curve
follow the commanded path within about a centimetre sideways on average under the cascaded geometric
controller. And the time to compute one action, $18$\,ms at
the operating point of D3IL-avoiding and $266$ to $275$\,ms on D3IL-aligning on the node of record,
prices the planner alone; the execution adds its own, on the s-curve's unprojected pair about
$5$\,ms per control step for the cascaded geometric controller and about $125$\,ms for MuJoCo MPC,
the controller and the simulator step together in each case and each in its own execution
environment, and the manipulator's inverse kinematics and joint controller were not timed. \guard{No loop of this thesis was run in real time;
the simulation waits for the planner (\autoref{sec:method:geometric}).}

One candidate plan is a tested option. On UAV-corridor, instantaneous-velocity matching at $\nfe = 20$
under endpoint projection over the hump costs $291$\,ms per action with one candidate against $422$
to $425$\,ms with four, and leaves $0.2$ violating steps against $0.4$ to $0.7$. With a single
candidate the three selection rules select the same plan, so success and control steps are identical
in all fifteen seed--geometry cells of every model on D3IL-avoiding; that is a fact about the rules,
not about the outcome of choosing among four, which is an outcome--cost decision per configuration.
With four candidates the generation time changes by at most six per cent, because the four are one
batched network evaluation, and the projection time grows by a factor of $3.2$ to $4.9$, because the
four programs are solved one after another on the CPU. Where four candidates are kept, the rule
matters by model: instantaneous-velocity matching holds $1.000$ under all three rules, while the two
average-velocity models are operated with temporal consistency, since cumulative projection cost, the
rule of \ac{DPCC}, occasionally selects a plan that stalls and costs them control steps, $72.4$ and
$97.2$ against $58.6$ and $59.4$ for analytic average-velocity matching at one and two evaluations.
In the candidate study of \autoref{sec:res:ablations}, where analytic average-velocity matching at one
evaluation measured $18.1$\,ms per control step, solving the four programs in parallel is estimated
at about $11.7$\,ms, an idealised prediction from two measured points.
\dataref{\autoref{sec:res:ablations}; \autoref{tab:avoiding-dpcc-protocol};
\autoref{app:uav-corridor-rules}.}

\section{Future Work}\label{sec:disc:future}

\begin{description}
  \item[Demonstrations for the quadrotor.] The demonstrations of UAV-corridor and UAV-s-curve are
    generated, one path forward from every position, and the s-curve's plans cut the corner that the
    demonstrated route clears. Human demonstrations of the quadrotor scenes with varied routes, from a
    pilot in the loop of a simulator or from flight logs, would test whether that diversity changes
    the relative performance of the objectives and the learned plans at the corner; NVIDIA Isaac Sim,
    a simulator with a rendered scene and a pilot interface, is the way to record them without a
    vehicle.
  \item[The controller and the commanded path.] Two things are open, and they differ. On the corridor
    the hand-over residue comes from the lead of the commanded position, $0.49$ to $0.61$\,m ahead of
    the vehicle at the first violating step of the projected flow-based flights; a plan that carries a
    velocity setpoint, or a controller that reads the next positions of the plan as feed-forward, is
    the tracking test for it. On the s-curve the commanded path itself cuts the corner the
    demonstrated route clears, and the lead does not explain the violations; the projector that binds
    the commanded position as well as the measured one, which the corridor uses and the s-curve did
    not fly, and the feasibility of the learned plans at the corner are the tests there.
  \item[Real-time execution.] The execution stack has to be timed and the loop run without waiting for
    the planner: the planner at one evaluation, the projector solved in parallel over the candidates,
    or with a single candidate where the rule selects the same plan, and the controller at its own rate.
  \item[The planning horizon.] Every plan of this thesis is eight steps long, and one step of it is
    executed before the next replan (\autoref{sec:setup:protocol:eval}); a constraint the horizon does
    not yet reach enters the projection at the replan that reaches it. A longer horizon would hand the
    projector the constraints that are in the program at every replan earlier; the span-switched ones --- the walls
    and the hump's two halves, \eqref{eq:method:env:switched} --- enter by the vehicle's current position, so they
    would need a rule on the predicted position, at the cost of an irrelevant wall over the whole horizon --- and
    the endpoint solve a
    farther endpoint, at a price: the
    network runs over a longer sequence, and the projection program carries one block of variables and
    one inequality per obstacle for every added step. Whether that foresight buys fewer violating
    steps or shorter routes at the added cost is untested.
  \item[Constraints from perception.] The projector consumes a short list of primitives with metric
    extent, halfspaces, keep-out disks and spheres and a workspace box, each in the planner's frame and
    each with its feasible side. A perception front-end that emits that list from depth or range
    sensing, and a map that keeps it beyond the field of view, is what a deployment outside the
    simulator needs; the uncertainty of the detector belongs in the inflation margin the projector
    already carries. The first experiment needs no detector: perturb the geometry the projector
    receives, by an offset, a radius error, a missing or a spurious primitive, and record where success
    with constraint satisfaction breaks.
  \item[Re-tasking at deployment.] Because the constraints never enter training, the same checkpoint
    can be asked to satisfy a new set: a closed passage, a keep-out volume, a tighter clearance. What
    has to be measured is the envelope, the slack a new set may leave before the projector pulls the
    plan off the data and success collapses rather than degrades, and the transfer of one checkpoint to
    a geometry its demonstrations never contained.
  \item[The alignment target.] The optional terminal cost of endpoint projection
    (\autoref{sec:method:hardflow}), zero in this thesis, is where the target of D3IL-aligning could
    enter the projection: the test is whether a task-aware terminal cost, one that represents the
    contact task and not the end-effector position alone, improves the box's progress to its target
    while keeping the reported constraint outcomes.
  \item[The sampler of instantaneous-velocity matching.] Sampling it from the scale it was trained at
    needs no retraining and would test the one sampling difference identified between it and the
    average-velocity models on D3IL-aligning and the quadrotor scenes.  \item[Seeds.] More training seeds on D3IL-aligning, UAV-corridor and UAV-s-curve.
\end{description}
\srcnote{Outlook material: ../future\_work/OUTLOOK\_20260910\_perception\_to\_constraints.md (the
projector's interface, the perturbation experiment) and
OUTLOOK\_20260910\_zero\_shot\_constraint\_retasking.md (constraints never enter training, verified in
code; the envelope experiment); the author's v4/notes.txt (expert data for the quadrotor; Isaac Sim named by the author,
v4.1); the planning horizon (author, v4.2; $H=8$, one executed step per replan and the per-step
constraint blocks: \autoref{sec:method:formalisation}, \autoref{sec:setup:protocol:eval}). Nothing here is a result.}
